\documentclass[ml]{sdoc}

\usepackage{biblatex}
\usepackage{changepage}

\begin{filecontents}[overwrite]{\jobname.bib}
    @book{bb,
        author = {B\'{e}la Bajnok},
        title = {Additive Combinatorics: A Menu of Research Problems},
        year = {2018},
        publisher = {CRC Press}
        }
        @misc{peter,
        author = {Peter Francis},
        title = {Personal Communication},
        URL = {https://bit.ly/3gJpzVf}
        }
        @article{lev,
        author = {Vsevolod Lev and Raphael Yuster},
        title = {On the Size of Dissociated Bases},
        year = {2011},
        journal = {The Electronic Journal of Combinatorics},
        volume = {18},
        DOI = {https://doi.org/10.37236/604}
        }
\end{filecontents}

\bibliography{\jobname}

\newcommand{\B}{\mathcal{B}}

\title{Maximal Dissociated Subsets of Finite Abelian Groups}
\author{Benjamin T. Shepard}
\dept{Department of Mathematics, Gettysburg College}
\email{shepbe01@gettysburg.edu}
\date{May 12, 2021}

\begin{document}

\maketitle

\begin{abstract}
    A subset $A$ of an abelian group $G$ is dissociated if $0\in G$ cannot be written as the signed sum of any of the elements of $A$. Furthermore, the dissociativity dimension of $A$, denoted $\dim A$, is the maximum size of a dissociated subset of $A$. We are interested in evaluating the function $\dim(G, m)$, which yields the minimum value of $\dim A$ when $|A|=m$. Here, we provide some basic values and bounds for this function, and evaluate it for certain elementary abelian $p$-groups. In particular, we obtain exact values when the lower bound is sharp for $p=2$ and $p=3$. We also disprove a conjecture made by B. Bajnok in 2018 regarding the exact value for cyclic groups.
\end{abstract}

\tableofcontents

\section{Introduction}

In this paper, when not specified, we will refer to $G$ as an arbitrary finite abelian group of order $n$ and exponent $\kappa$, and take $m$ to be a positive integer. We begin with some definitions.

\begin{definition}
    A subset $A=\{a_1,\ldots,a_m\}\subseteq G$ is dissociated if every possible equality of the form
    \[\sum_{i=1}^m \lambda_i a_i = 0,\]
    where $\lambda_i\in\{-1,0,1\}$, implies that $\lambda_i=0$ for all $i$. We usually denote the set of all possible sums as defined above by $\Sigma A$, and write $0\notin\Sigma A$ if $A$ is dissociated.
\end{definition}

Dissociated subsets of a group have been investigated, however there is an interest in finding dissociated subsets of any subset of a group. For this, we define

\begin{definition}
Let $A\subseteq G$. The quantity
    \[\dim A \defeq \max\{|D|\mid D \subseteq A,\, D\ \mathrm{is\ dissociated}\}\]
    is known as the dissociativity dimension of $A$ in $G$.
\end{definition}

Throughout this paper, we refer to $\dim A$ as simply the dimension of $A$. Finally, we will define the main function that we are interested in investigating.

\begin{definition}
Given $G$ and $m\leq n$, define
    \[\dim(G,m)\defeq\min \{\dim A\mid A \subseteq G,\, |A|=m\};\]
    that is, the minimum dissociativity dimension of any subset of $G$ of size $m$.
\end{definition}

\section{Previous Results}

We currently have the following:

\begin{theorem}[Lev and Luster, 2011; \cite{lev}]\label{F.188}
    For all $G$ and $A\subseteq G$, we have
    \[r_A\leq \dim A\leq \lfloor r_A\cdot\log_2\,\kappa\rfloor\]
    where $r_A$ is the rank of the subgroup $\langle A\rangle$ generated by $A$.
\end{theorem}

\begin{proposition}[Bajnok, 2018; \cite{bb}]\label{nr}
    If $G$ has order $n$ and invariant factorization
    \[G\cong\Z_{n_1}\times\cdots\times\Z_{n_r},\]
    then
    \[\lfloor\log_2 n_1\rfloor+\ldots+\lfloor\log_2 n_r\rfloor\leq\dim(G,n)\leq\lfloor\log_2 n\rfloor.\]
    In particular, for all positive integers $n$, we have
    \[\dim(\Z_n,n)=\lfloor\log_2 n\rfloor.\]
\end{proposition}

\begin{proposition}[Bajnok, 2018; \cite{bb}]\label{123}
    For all $G$, we have $\dim(G, 1)=0$, $\dim(G, 2)=1$, and
    \[\dim(G, 3) = \begin{cases}1 & \text{if $\kappa\geq3$;} \\ 2 & \text{if $\kappa=2$.}\end{cases}\]
\end{proposition}

\section{Main Results}

\subsection*{Elementary Results}
\addcontentsline{toc}{subsection}{Elementary Results}

We first evaluate the cases of $m=4$ and $m=5$ for all $G$ to obtain an extension of Proposition \ref{123}.

\begin{proposition}\label{45}
    For all $G$, we have $\dim(G,4)=2$ and
    \[\dim(G,5)=\begin{cases}
    2 & \text{if $\kappa\geq3$;}\\
    3 & \text{if $\kappa=2$.}
    \end{cases}\]
\end{proposition}
\begin{proof}
    When $\kappa=2$, both claims follow from Theorem \ref{Z2r}. Therefore, assume $\kappa\geq3$. 
    For $m=4$, let
    \[A = \{0, g_1, -g_1, g_2\}\]
    such that $g_1, -g_1$, and $g_2$ are all nonzero and distinct. For $m=5$, a non-boolean abelian group $G$ of order at least 5 contains at least two elements $g_1,g_2\in G$ such that $g_1,g_2\neq 0,-g_1,-g_2$, so we can let
    \[A_1= A\cup\{-g_2\}\]
    such that $A\cap\{-g_2\}=\emptyset$. Clearly both $A$ and $A_1$ do not have a dissociated subset of size 3, so $\dim A,A_1\geq 2$. In both $A$ and $A_1$, one can easily see that $0\notin\Sigma\{g_1,g_2\}$, so $\dim A =\dim A_1=2$. Now, since only subsets of a set of the form $\{0,g_1,-g_1\}$ have dimension less than 2, we have $\dim(G,4)=\dim(G,5)=2$ as claimed.
\end{proof}

% \begin{corollary}
%     For all $G$ we have
%     \begin{align*}
%         \dim(G,1)&=0\\
%         \dim(G,2)&=1\\[5pt]
%         \dim(G,3)&=\begin{cases}
%             1 & \text{if $\kappa\geq3$;} \\ 
%             2 & \text{if $\kappa=2$.}
%             \end{cases}\\[5pt]
%         \dim(G,4)&=2\\[5pt]
%         \dim(G,5)&=\begin{cases}
%             2 & \text{if $\kappa\geq3$;}\\
%             3 & \text{if $\kappa=2$.}
%             \end{cases}
%     \end{align*}
% \end{corollary}

We next introduce several bounds. This first proposition states that the dimension function is monotonically increasing over $m$, which is an important fact that will be used later.

\begin{proposition}\label{m1m2}
    If $m_1$ and $m_2$ are positive integers such that $m_1\geq m_2$, then
    \[\dim(G,m_1)\geq\dim(G,m_2).\]
\end{proposition}
\begin{proof}
    For simplicity, let $\delta_1=\dim(G, m_1)$ and $\delta_2=\dim(G, m_2)$. We must prove that every subset of $G$ of size $m_1$ has dimension at least $\delta_2$. Let $A\subseteq G$ be arbitrary and of size $m_1$.
    Since $|A|\geq m_2$, $A$ contains at least one subset $B$ of size $m_2$. By definition of $\delta_2$, every subset of $G$ of size $m_2$ has dimension at least $\delta_2$, so $\dim B\geq \delta_2$. Thus $B$ contains a dissociated subset $D$ of size $\delta_2$, and $B\subseteq A$, so we have $D\subseteq A$. By definition of dimension, $\dim A\geq |D|=\delta_2$ as claimed.
\end{proof}

\begin{proposition}\label{logn}
    For all $G$ of order $n$, and all $m\leq n$, we have
    \[\dim(G,m)\leq\lfloor \log_2 n \rfloor.\]
\end{proposition}
\begin{proof}
    By Proposition \ref{nr}, $\dim(G,n)\leq\lfloor\log_2 n\rfloor.$ Using Proposition \ref{m1m2}, we obtain
    \[\dim(G,m)\leq\dim(G,n)\leq\lfloor\log_2 m\rfloor\]
    as claimed.
\end{proof}

\begin{proposition}\label{m/2}
    Let $G$ have order $n$. For $G$ for which $|\!\operatorname{Ord}(G,2)|<m/2$, we have \[\dim(G,m)\leq \left\lfloor \frac{m}{2}\right\rfloor.\]
\end{proposition}
\begin{proof}
    Let $k=\lfloor m/2\rfloor$. If $m$ is even, then $m=2k$; if $m$ is odd, then $m=2k+1$. Since less than $m/2$ elements in $G$ have order 2, we may let
    \[A=\{0\}\cup\bigcup_{i=1}^k \{a_i,-a_i\}\subset G\]
    such that $a_i\neq -a_i\neq 0$ for all $i$ and
    \[\{a_i,-a_i\}\cap\{a_j,-a_j\}=\emptyset\]
    for all $i\neq j$. If $m$ is even, take $A$ without $-a_k$. In either case, $|A|=m$. 
    
    Suppose that there exists a dissociated subset $D\subset A$ such that $|D|> k$. Since the smallest possible size of $D$ is $k+1$, if we choose $k+1$ elements of $A$, we are forced to choose either 0 or both $a_i$ and $-a_i$ for some $i\in \{1,\ldots,k\}$. In either case, $0\in \Sigma D$, so we must have $|D|\leq k$, a contradiction. Therefore, $\dim A\leq k$, which directly implies that $\dim(G,m)\leq k$, as claimed.
\end{proof}

Note that these bounds are not generally sharp. For large $m$, Proposition \ref{logn} is better, but note that it is in terms of $n$, which loosens the bound.

\subsection*{Lower Bounds}
\addcontentsline{toc}{subsection}{Lower Bounds}

In the next theorem, we introduce a lower bound for all $G$ and $m$. First, we have a useful lemma, whose proof is omitted here.

\begin{lemma}\label{subgroup}
    If $G$ has invariant factorization
    \[G\cong\Z_{n_1}\times\cdots\times\Z_{n_r},\]
    then the group
    \[\Z_{n_{r-k+1}}\times\cdots\times\Z_{n_r}\]
    is the largest subgroup of $G$ with rank $k$.
\end{lemma}

This allows us to obtain the following result:

\begin{theorem}\label{lower}
    Suppose that $G$ has invariant factorization
    \[G\cong\Z_{n_1}\times\cdots\times\Z_{n_r}.\]
    Then for all $m$ we have
    \[\dim(G, m)\geq k\]
    for the unique $k$ such that 
    \[\prod_{i=r}^{r-k+2}\!n_i<m\leq\prod_{i=r}^{r-k+1}\!n_i.\]
\end{theorem}
\begin{proof}
    Let $A\subseteq G$ have size $m$. Since
    \[\prod_{i=r}^{r-k+2}\!n_i<m,\]
    the size of $A$ is strictly greater than that of the subgroup
    \[\Z_{n_{r-k+2}}\times\cdots\times\Z_{n_r}.\]
    By Lemma \ref{subgroup}, this is the largest subgroup of $G$ with rank $k-1$. Therefore, $\langle A \rangle$ has rank $r_A\geq k$. Now from Theorem \ref{F.188} we have $\dim A\geq k$, and the claim follows.
\end{proof}

From this we immediately obtain

\begin{corollary}\label{Zkr}
    For all positive integers $k$ and $r$, we have
    \[\dim(\Z_k^r,m)\geq \lceil \log_k m\rceil.\]
\end{corollary}
\begin{proof}
    Since $k^{i-1}<m\leq k^i$ for some $i\in\{1,\ldots,r\}$, we have $i\geq\lceil\log_k m\rceil$. The claim now follows from Theorem \ref{lower}.
\end{proof}

It should be noted that for $k>3$, equality does not generally hold. For example, it is an immediate consequence from Theorem \ref{F.188} that
\[\dim(\Z_k^r,k^r)\geq kr\]
for all $r$ and $k\geq4$, which is larger than the lower bound of $r$ given here. We do not currently know much about other values of $m$.

\subsection*{Exact Values}
\addcontentsline{toc}{subsection}{Exact Values}

It turns out that the bound in Corollary \ref{Zkr} is sharp for certain elementary abelian $p$-groups; in particular, $p=2$ and $p=3$.

% The following theorems focus on elementary abelian $p$-groups $\Z_p^r$, for which we can form a lower bound and obtain some exact values.
% \begin{theorem}\label{Zp}
% Let $p$ be a prime. For all $p$, and all positive integers $r$ and $m\leq p^r$, $$\dim(\Z_p^r,m)\geq\lceil\log_p m\rceil.$$
% \end{theorem}
% \begin{proof}
%     Let $A\subseteq\Z_p^r$ with $|A|=m$. Since $p$ is prime, the subgroups of $\Z_p^r$ are of the form $\Z_p^{a}$ for positive integers $a\leq p$. Thus, $\langle A \rangle\cong\Z_p^k$ for some $k\leq p$. Then $|\langle A \rangle|=p^k$, and since $|A|\leq|\langle A\rangle|$, we have
%     \begin{align*}
%         |A|&\leq p^k\\
%         m&\leq p^k\\
%         \log_p m&\leq k\\
%         \lceil \log_p m \rceil&\leq k
%     \end{align*}
% where the last inequality holds since $k$ is an integer. This gives us a lower bound on the rank of $\langle A \rangle$, so by Proposition \ref{F.188}, we have $\dim A\geq \lceil \log_p m \rceil$. Since $A$ is arbitrary, $\dim(\Z_p^r,m)\geq \lceil \log_p m \rceil$, as claimed.
% \end{proof}

% It should be noted that for $p>3$, equality does not generally hold. We do not currently know much about these cases. However, it turns out that this bound is sharp for $p=2$ and $p=3$. We obtain both

\begin{theorem}\label{Z2r}
    For all positive integers $r$ and $m\leq 2^r$, we have \[\dim(\Z_2^r,m)=\lceil\log_2 m\rceil.\]
\end{theorem}
\begin{theorem}\label{Z3r}
    For all positive integers $r$ and $m\leq3^r$, we have
    \[\dim(\Z_3^r,m)=\lceil \log_3 m \rceil.\]
\end{theorem}

The proofs of both theorems are similar, so here we prove both using a general construction.
\begin{proof}
    Let $p\in\{2,3\}$. Since $p$ is prime, $\Z_p$ is a field and $\Z_p^r$ forms a vector space over $\Z_p$. Recall that for any $A\subseteq \Z_p^r$, the coefficients of sums in $\Sigma A$ are given by $\lambda_i=\{0,1,-1\}$. The scalars in $\Z_3^r$ are exactly $\lambda_i$. Also, since each element of $\Z_2^r$ is its own inverse, for $p=2$ we only need to consider $\lambda_i=\{0,1\}$, which are exactly the scalars in $\Z_2^r$. Therefore, for any $X\subseteq \Z_p^r$, we have an equivalence: $X$ is dissociated if and only if the vectors in $X$ are linearly independent. 
    
    We now use Proposition \ref{m1m2} to narrow down the cases; it suffices to prove that $\dim(\Z_p^r,p^k)=k$ and $\dim(\Z_p^r,p^k+1)=k+1$ for all positive integers $k\leq r$.
    
    First, let $m=p^k$. From Theorem \ref{Zkr}, we have $\dim(\Z_p^r,p^k)\geq k$, so it suffices to show the existence of subset of $\Z_p^r$ of size $p^k$ and dimension $k$. Let
    \[A=\{0\}^{r-k}\times\Z_p^k.\]
    It is clear to see that $\dim A=\dim\Z_p^k$. Now, it is a well known fact that if a vector space $V$ is spanned by $k$ vectors, and $l>k$, then any set of $l$ vectors in $V$ is linearly dependent. Since the standard basis for $\Z_p^k$,
    \[\B=\{(e_1,\ldots,e_k)\mid e_i=1 \text{ for some }i\in \{1,\ldots,k\},\,e_j=0\text{ for }j\neq i\},\]
    spans $\Z_p^k$ and has size $k$, any subset of $\Z_p^k$ must have size at most $k$ to be dissociated. It is clear that $0\notin\Sigma\B$ since $\B$ is linearly independent, so $\dim\Z_p^k=k$. Therefore, we obtain $\dim A = k$ as desired.
    
    \vspace{1em}
    Now let $m=p^k+1$. Theorem \ref{Zkr} yields $\dim(\Z_p^r, p^k+1)\geq k+1$, so it again suffices to show existence of a subset of $\Z_p^r$ of size $p^k+1$ with dimension $k+1$. Write
    \[A_0=A\cup\B_0\]
    where
    \[\B_0=\{1\}\times \{0\}^{r-1}.\]
    Since $\B\cap\B_0=\emptyset$, the set
    \[(\{0\}^{r-k}\times\B)\cup\B_0 \subset A_0\]
    has size $k+1$ and is clearly dissociated. 
    
    Now, suppose that there exists a dissociated subset $D\subset A_0$ with $|D|> k+1$. Since $\dim A =k$,  $D$ cannot contain more than $k$ elements of $A$. But since we have $|A_0|=k+1$, it is impossible to put two more elements of $A_0$ into $D$. Hence $|D|\leq k+1$, a contradiction. From this, we obtain $\dim A_0=k+1$ as desired.
\end{proof}

This proof does not work for $p\geq 5$; even though $\Z_p^r$ is a vector space for all values of $p$, the only groups where equivalence holds between dissociativity and linear independence are $\Z_2^r$ and $\Z_3^r$.

Using Proposition \ref{nr}, we can easily get equality for a few related groups when $m$ is maximal.

\begin{proposition}
    For all positive integers $r$ and $k$, we have
    \[\dim(\Z_2^r\times\Z_{2k},2^{r+1}k)=\lfloor\log_2 2^{r+1}k\rfloor.\]
\end{proposition}
\begin{proof}
    From Proposition \ref{nr}, we have
    \[r+\lfloor\log_2 2k\rfloor\leq\dim(\Z_2^r\times\Z_{2k},k2^{r+1})\leq\lfloor\log_2 k2^{r+1}\rfloor.\]
    Since clearly
    \[r+\lfloor\log_2 2k\rfloor=r+1+\lfloor\log_2 k\rfloor\]
    and
    \[\lfloor\log_2 k2^{r+1}\rfloor=\lfloor\log_2 2^r\cdot 2\rfloor+\lfloor\log_2 k\rfloor=r+1+\lfloor\log_2 k\rfloor,\]
    the bounds are equal and equality holds.
\end{proof}

\subsection*{Cyclic Groups}
\addcontentsline{toc}{subsection}{Cyclic Groups}

We now move on to cyclic groups $\Z_n$. In \cite{bb}, B. Bajnok conjectured that for all positive integers $n$ and $m\leq n$, the following equality holds:
\[\dim(\Z_n,m)=\lfloor \log_2 m\rfloor.\]

Using the code provided by Francis in \cite{peter}, we have the following table of values that disagree with this conjecture:

\vspace{1em}
\begin{table}[H]
\centering
\begin{tabular}{c|c|c|c}
$G$  & $m$ & $\dim(G, m)$ & $\lfloor \log_2 m\rfloor$\\ \hline
$\Z_{17}$ & $14,15$ & 4   & 3 \\ \hline
$\Z_{19}$ & $14,15$ & 4   & 3        \\ \hline
$\Z_{22}$ & 15      & 4   & 3      \\ \hline
$\Z_{23}$ & $14,15$ & 4   & 3      \\ \hline
$\Z_{26}$ & 15      & 4     & 3     \\ \hline
$\Z_{28}$ & 15      & 4    & 3    \\ \hline
$\Z_{29}$ & $14,15$ & 4   & 3           
\end{tabular}
\caption{Values larger than the conjecture.}
\end{table}
\vspace{1em}

All other groups below $n=30$ agree exactly for all values of $m$. In light of these results, I believe that the lower bound still holds. 

I propose the following more modest conjecture:

\begin{conjecture}\label{newZn}
    For all positive integers $n$ and $m\leq n$, we have
    \[\dim(\Z_n,m)\geq\lfloor\log_2 m\rfloor.\]
\end{conjecture}

I also believe that powers of two agree with Bajnok's conjecture.

\begin{conjecture}\label{newZn2^k}
    For all positive integers $n$ and $m\leq n$, we have
    \[\dim(\Z_n,2^k)=k.\]
\end{conjecture}

It may also be more feasible to consider specific cyclic groups. For instance, the case when $n$ is a power of two may be of particular interest.

\section{Future Work}

Future projects should explore the validity of Conjectures \ref{newZn} and \ref{newZn2^k}. Additionally, since the upper bounds of Proposition \ref{m/2} and \ref{logn} are not always sharp, future work should look into the possibility of improving these bounds. 

Furthermore, since Corollary \ref{Zkr} is only sharp for $k=2$ and $k=3$, future work should attempt to improve the lower bound for $k\geq4$, and possibly other elementary abelian $p$-groups for $p\geq5$.\vskip1em\relax

{\bf Acknowledgments.} I would like to thank Dr. Bela Bajnok for his guidance throughout the duration of this research. I would also like to thank Peter Francis, for providing the code used to generate the cyclic group table (among many others), and Matt Torrence, for his general help throughout the semester. Finally, I would like to thank Dr. Michael Shepard for his contribution to a separate program for evaluating cyclic groups.

% \nocite{*}
\printbibliography

\end{document}
